% =====================================================================
%  chunks/rev_ch01.tex —— 第 1 章 读书笔记（Lead 撰）
% =====================================================================
\chapterreview{%
第 1 章篇幅不过十页，却是全书的方法论总纲。作者刻意不引入任何形式定义，只挑了五个彼此相距极远的例子：Euler 的 $V-E+F=2$、三管道三水井问题、Gauss 曲率的积分、Jordan 曲线定理与环绕系数、以及复分析里的 Cauchy 积分与单值性。这五个例子的价值不在结论，而在于它们各自示范了一种「把几何问题翻译成不变量语言」的路数，且五条路互不相同：组合的（重分不变性）、图论的（嵌入障碍）、微分几何的（曲率积分）、代数拓扑的（环绕系数、基本群）、分析学的（积分在同伦下的不变性）。

最值得注意的是同一不变量的多重面目。$2$ 这个数在 Euler、Maxwell（山峰加坑洼减垭口）、Gauss--Bonnet 三处反复出现，作者用「漂浮海岛」把 Maxwell 的观察强行补齐成闭曲面上的断言，正是要说明：三条定理是同一条拓扑性质的三种坐标系。等到亏格 $g$ 出现，$2$ 变成 $2-2g$，三者同时被修正，这条线索才算收束。读到这里应当意识到，后面全书的示性类、指标定理，本质上都是在为「同一个不变量的多种表述」建立统一的机器。

第二条主线是基本群的登场。复变函数的多值性被翻译成基本群到对称群的同态（单值性），这条线索把一个纯分析的现象变成了群论对象；纽结补集的基本群是否交换，则判定纽结是否平凡。作者由此点明：非交换性是这两类问题真正的难点，而交换化之后得到的 $H_1$ 恰恰丢掉了关键信息。这为第 3 章为什么要发展同调论、又为什么要回到非单连通空间埋下伏笔。

第三处伏笔是末尾的 Fredholm 算子与指标。奇异积分方程的指标是同伦不变量一句，已经把分析问题彻底拓扑化；Atiyah--Singer 指标定理在此只是被点名，未展开。读到第 4 章 §2 的 (2.26) 与 §5 再回看这一页，会明白作者为什么把这个例子放在第一章的收尾——它预示了「指标 = 拓扑不变量」这条贯穿全书的主线。

若要用一句话概括本章的读法：\emph{不要去找定义，要找「同一个量在不同语言里的化身」}。作者列的五个例子，每一个都至少给出两种化身；全书后续每一节几乎都在重复这个动作。第 1 章还顺手立下了本书的两个表述习惯：拓扑性质指在同伦（连续形变）下稳定的性质；而「稳定」的严格化，将分别由同调群、示性类、配边群来承担。
}

% CHUNK rev_ch01 DONE
